\documentclass[11pt,letterpaper]{article}

\usepackage[T1]{fontenc}
\usepackage{amsmath,amssymb}
\usepackage{tikz}
\usepackage{pgfplots}
\pgfplotsset{compat=1.18}
\usepackage{algorithm}
\usepackage{algpseudocode}
\usepackage{booktabs}
\usepackage{enumitem}
\usepackage[hidelinks]{hyperref}

\newcommand{\N}{\mathbb{N}}
\newcommand{\Z}{\mathbb{Z}}
\newcommand{\Q}{\mathbb{Q}}
\newcommand{\R}{\mathbb{R}}
\newcommand{\C}{\mathbb{C}}
\newcommand{\Map}[3]{#1 \colon #2 \to #3}
\newcommand{\Mod}[1]{#1\text{-}\mathrm{Mod}}
\newcommand{\GroupHom}{\mathrm{Hom}_{\mathbf{Grp}}}
\newcommand{\RingHom}{\mathrm{Hom}_{\mathbf{Ring}}}

\newtheorem{theorem}{Theorem}
\newtheorem{lemma}[theorem]{Lemma}
\newtheorem{corollary}[theorem]{Corollary}
\newtheorem{definition}[theorem]{Definition}
\newtheorem{remark}[theorem]{Remark}
\newenvironment{proof}{\noindent\textit{Proof.}}{\hfill\(\square\)}

\title{Positive Two-Affine Closures of Even Length}
\author{}
\date{\today}

\begin{document}
\maketitle

\begin{abstract}
Every positive two-affine closure of even shortcut length has equal letters and unit components. The proof orders the letters and combines the closure determinant with its difference-of-squares factorization. The unit family is realized for every positive half length; in particular no unequal positive components occur at even length.
\end{abstract}

\section{Two-affine closure}

Fix a positive integer multiplier \(q\) and positive integers \(a\) and \(b\). A two-affine word on components \(n\) and \(y\) is the pair of equations
\begin{equation}
\label{eq:two-affine}
q n + 1 = 2^{a} y, \qquad q y + 1 = 2^{b} n.
\end{equation}
The length is \(K = a + b\). The equal-component relation is the special case
\begin{equation}
\label{eq:equal-closure}
q n + 1 = 2^{a} n.
\end{equation}

\begin{lemma}
\label{lem:letters}
If \(n > 0\) and both \(q n + 1 = 2^{a} n\) and \(q n + 1 = 2^{b} n\), then \(a = b\).
\end{lemma}

\begin{proof}
The two right-hand sides are equal, so \(2^{a} n = 2^{b} n\). Since \(n > 0\), cancelling \(n\) gives \(2^{a} = 2^{b}\). It remains to recover the exponent. Proceed by induction on \(a\).

If \(a = 0\), then \(1 = 2^{b}\). The case \(b = 0\) is immediate. If \(b = j + 1\), then \(1 = 2^{j} \cdot 2\), so reduction modulo \(2\) gives \(1 = 0\), which is impossible.

If \(a = k + 1\) and \(b = 0\), the same reduction gives \(0 = 1\). If \(b = j + 1\), the equation is \(2 \cdot 2^{k} = 2 \cdot 2^{j}\). Dividing by \(2\) yields \(2^{k} = 2^{j}\), so the inductive hypothesis gives \(k = j\), hence \(a = b\).
\end{proof}

\begin{lemma}
\label{lem:even-length}
For every natural number \(a\), the sum \(a + a\) is even. Consequently, if \(n = y > 0\) solves \eqref{eq:two-affine}, then \(a = b\) and \(K\) is even.
\end{lemma}

\begin{proof}
The sum \(a + a\) is twice \(a\). Substituting \(y = n\) into \eqref{eq:two-affine} produces the two equal-component equations of Lemma~\ref{lem:letters}, so \(a = b\) and \(K = a + a\).
\end{proof}

\begin{lemma}
\label{lem:cofactor}
If \(n > 0\) and \(q n + 1 = 2^{a} n\), then \((2^{a} - q) n = 1\). In particular \(n = 1\) and \(2^{a} = q + 1\).
\end{lemma}

\begin{proof}
The hypothesis rearranges as \(q n < 2^{a} n\), so
\[
2^{a} n - q n = 1.
\]
The left side is \((2^{a} - q) n\). Thus the product of two natural numbers is \(1\). Neither factor can be \(0\), because \(n > 0\) and the product is positive. If the first factor is at least \(2\), then the product is at least \(2\). Both factors are therefore \(1\), which means \(n = 1\) and \(2^{a} - q = 1\). Adding \(q\) gives \(2^{a} = q + 1\).
\end{proof}

\begin{lemma}
\label{lem:squares}
If \(q \le A\), then \((A - q)(A + q) = A^{2} - q^{2}\). Taking \(A = 2^{m}\) gives
\[
(2^{m} - q)(2^{m} + q) = 2^{2m} - q^{2}.
\]
\end{lemma}

\begin{proof}
Expand the product as \((A - q)A + (A - q)q\). The inequalities \(q A \le A^{2}\) and \(q^{2} \le A q\) make the truncated subtractions exact, and the intermediate term \(A q\) cancels, leaving \(A^{2} - q^{2}\). The power form is the case \(A = 2^{m}\), using \(2^{m} \cdot 2^{m} = 2^{2m}\).
\end{proof}

\begin{lemma}
\label{lem:seed}
If \(n > 0\) and \(q n + 1 = 2^{a} n\), then
\[
n(2^{2a} - q^{2}) = q + 2^{a}.
\]
Conversely, if \(q \le 2^{a}\) and \(n(2^{2a} - q^{2}) = q + 2^{a}\), then \((2^{a} - q) n = 1\).
\end{lemma}

\begin{proof}
Lemma~\ref{lem:cofactor} gives \((2^{a} - q) n = 1\) and \(q < 2^{a}\). Lemma~\ref{lem:squares} factors the determinant. Multiplying the cofactor equation by \(2^{a} + q\) produces the determinant equation. In the other direction, factor \(2^{2a} - q^{2}\) and cancel the positive factor \(2^{a} + q\).
\end{proof}

\begin{lemma}
\label{lem:odd-multiplier}
If \(a \ge 1\), \(n\) is odd, and \(q n + 1 = 2^{a} n\), then \(q\) is odd.
\end{lemma}

\begin{proof}
Since \(a \ge 1\), one has \(2^{a} = 2^{a-1} + 2^{a-1}\), so \(2^{a}\) is even and \(2^{a} n\) is even. The left side \(q n + 1\) is therefore even, which means \(q n\) is odd. A product is odd only when both factors are odd.
\end{proof}

\begin{theorem}
\label{thm:unit}
For every \(a \ge 1\), the multiplier \(q = 2^{a} - 1\) and the state \(n = 1\) satisfy \(q n + 1 = 2^{a} n\). The multiplier is odd, and the state does not satisfy \(n \ge 3\).
\end{theorem}

\begin{proof}
Substitute \(n = 1\): \((2^{a} - 1) + 1 = 2^{a}\). For \(a \ge 1\), the power \(2^{a}\) is even and at least \(2\), so \(2^{a} - 1\) is odd. The inequality \(3 \le 1\) is false.
\end{proof}

\begin{theorem}
\label{thm:no-equal}
Let \(n = y\), and assume \(n \ge 3\) and \(y \ge 3\). Then \eqref{eq:two-affine} has no solution in positive integers \(q, a, b\). If in addition \(K\) is even, \(a \ge 1\), \(b \ge 1\), and \(n\) and \(y\) are odd, the same obstruction applies: the half length \(t = K/2\) equals both letters, the determinant equation of Lemma~\ref{lem:seed} holds at that half, the multiplier is odd, and the state is still the unit excluded by the lower bound \(3\).
\end{theorem}

\begin{proof}
By Lemma~\ref{lem:even-length}, \(a = b\) and \(K = 2a\). Write \(t\) for the common half, so \(a = b = t\). Lemma~\ref{lem:cofactor} gives \(n = 1\). This contradicts \(n \ge 3\). The same contradiction applies to \(y\), because \(y = n\).

Now assume \(K\) is even, and write \(K = t + t\) from the definition of evenness. The previous paragraph still forces \(a = b\), so \(a + a = t + t\), hence \(a = t\). The hypotheses \(a \ge 1\) and \(b \ge 1\) give \(2 \le K = t + t\), so \(t \ge 1\). Lemma~\ref{lem:odd-multiplier} makes \(q\) odd, Lemma~\ref{lem:seed} is the determinant equation at this half, and Lemma~\ref{lem:cofactor} again forces \(n = 1\), contradicting \(n \ge 3\).
\end{proof}

\begin{theorem}
\label{thm:odd-length}
If \(K\) is odd and \(n > 0\), there is no solution of \(q n + 1 = 2^{a} n = 2^{b} n\).
\end{theorem}

\begin{proof}
Lemma~\ref{lem:letters} forces \(a = b\), so \(K = a + a\) is even by Lemma~\ref{lem:even-length}. An odd length cannot be even.
\end{proof}

\section{All even lengths}

\begin{theorem}[Complete even-length classification]
\label{thm:even-complete}
Let \(q,a,b,n,y\) be positive integers satisfying \eqref{eq:two-affine}.
If \(a+b=2t\), then
\[
a=b=t,\qquad q=2^t-1,\qquad n=y=1.
\]
Conversely these values give a closure for every \(t\ge1\).
\end{theorem}

\begin{proof}
Interchange \((a,n)\) with \((b,y)\), if needed, so that \(a\le b\).
Then \(a\le t\).
Eliminating \(y\) gives
\[
\delta n=q+2^a,\qquad \delta=2^{2t}-q^2.
\]
The right side is positive and \(n>0\), so \(\delta>0\).
Put \(A=2^t\). Then \(q<A\), hence \(A-q\ge1\), and
\(\delta=(A-q)(A+q)\ge A+q\).
Since \(n\ge1\) and \(2^a\le A\),
\[
\delta\le\delta n=q+2^a\le q+A\le\delta.
\]
Every inequality is equality.
Thus \(n=1\), \(2^a=A\), and \(A-q=1\).
Injectivity of powers of two gives \(a=t\); then \(b=t\) and \(q=A-1\).
The first closure equation now gives \(Ay=A\), so \(y=1\).
Undoing the interchange gives the same conclusion.
The converse is direct substitution, and \(t\ge1\) makes \(q\) positive.

The active declarations \texttt{EvenClosureClassification} and
\texttt{UnitClosure} in \texttt{TwoAffineClassification.lean} formalize both directions.
The length \(a+b\) counts shortcut instructions; the repeated unit block is not asserted to be a primitive two-component cycle.
\end{proof}

\end{document}
